%!TEX root = ./../../thesis-trier.tex

\IN{completely reachable!automaton}
Bondar and Volkov~\cite{BondarV16,BondarCV23}
called $\semiautomat{A}{Q}{\Sigma}{\delta}$ \emph{completely reachable} if  every non-empty 
subset $S \subseteq Q$ is reachable, i.e., 
there exists a word $u \in \Sigma^*$ such that $\applytrans{\delta}{Q}{u} = S$. 
Equivalently, in the power automaton with the empty set removed, $Q$ is the unique minimum element for the reachability quasiorder.

Completely reachable automata are strongly connected and synchronizing.

\begin{lemma}[Bondar \& Volkov~\cite{BondarV16,BondarCV23}]
	\label{prelim:lem:compl_reach_implies_SCC_sync} 
	If $\semiautomat{A}{Q}{\Sigma}{\delta}$ is completely reachable, then~$\mathcal A$
	is synchronizing and strongly connected.
\end{lemma}
\begin{proof}
	Let $p, q \in Q$.
	Then $\applytrans{\delta}{Q}{u} = \{q\}$ for some word $u \in \Sigma^*$.
	Thus $u \in \Syn(\mathcal A)$ and since $\applytrans{\delta}{p}{u} = q$
	we also conclude that~$\mathcal A$ is strongly connected.
\end{proof}

In a completely reachable automaton there exists
for each non-empty subset a word mapping the whole state set to this subset.
So, its transformation monoid has size at least $2^n - 1$.

Next we show that the automata from the \Cerny family (\cref{intro:sec:cerny,intro:fig:cerny})
are completely reachable for each $n \ge 1$. This fact also follows from \cref{sync_max:thm:circular_sync_max}, a more general result
stated later, but it might be instructive to have a concrete example together with a self-contained proof.

\begin{theorem}
	\label{sync_max:thm:cerny_compl_reach}
	The automata $\mathcal C_n$ from the \Cerny family (Section~\ref{intro:sec:cerny} \& Figure~\ref{intro:fig:cerny})
	are completely reachable.
\end{theorem}
\begin{proof}
	Let $\semiautomat{C_n}{[n]}{\{a,b\}}{\delta_n}$.
	The whole state set is reachable by the empty word.
	Next, consider a non-empty and proper set $S \subseteq Q$.
	We have $0 \notin \delta_n(S, b^i)$ and $1 \in \delta_n(S, b^i)$
	for some $0 \le i < n$.
	Inductively, there exists $u \in \{a,b\}^*$ such that
	$\delta_n(S, b^i) \cup \{0\} = \delta_n(Q, u)$.
	Then we have $\delta(Q, uab^{n - i}) = S$.
%	
%	An alternative proof can be obtained by considering the words $u_i = b^{n - i} a b^i$ for $i \in [n]$
%	for which we have $\delta(j, u_i) = j + 1$ for $j = i$
%	and $i \in \{0,1,\ldots,n-2\}$, $\delta(n-1, u_{n-1}) = 0$,
%	and $\delta(j, u_i) = j$ in all other cases. 
%	Then the claims follows since for distinct $i_0, \ldots, i_m \in [n]$
%	with $\delta(i_j, b^{i_{j+1} - i_{j}}) = i_{j+1}$
%	for $j \in \{0,\ldots,n-2\}$
%	and $i_m + 1 \not\equiv i_1 \pmod{n}$
%	%and $i_j + s_j \equiv i_{j+1}
%	%with $0 \le i_1$ and $i_m \le n - 2$,
%	%or $1 \le i_1$ and  $i_m \le n - 1$,
%	we have
%	$
%		\delta([n], u_{i_1} \cdots u_{i_m}) = [n] \setminus \{ i_1, \ldots, i_m \}.
%	$
%	For example, for $n > 2$ we have $\delta([n], u_{n-1} u_0) = \{1,\ldots,n-2\}$,
%	but note that $\delta([n], u_0 u_{n-1}) = \{0,\ldots,n-2\}$.	
\end{proof}